\section{加密矿企正在转向 AI 数据中心}

自2024年以来，比特币减半周期、比特币价格回调及市场竞争加剧等因素共同作用，持续压缩加密矿企传统业务的利润空间。多家矿企宣布转向数据中心业务，希望将现有数字基础设施资产重新部署到快速增长的 AI 市场。

今年，数据中心转型进一步进入落地阶段。MARA 着手收购数据中心及已具备供电条件的场地，取得法国电力公司 EDF 旗下 Exaion 的控股权，并将数字基础设施确定为主要增长方向，计划逐步把现有矿场发展为 AI／高性能计算（HPC）园区。Riot 第二季度的数据中心收入达到2,320万美元，并与一家领先的前沿 AI 实验室签订为期20年的数据中心协议，在其 Rockdale 园区提供191兆瓦的关键 IT 容量。2026年2月，Cipher 从 Cipher Mining 更名为 Cipher Digital，将业务定位扩展至数据中心开发与运营。

在挖矿经济性承压的背景下，加密矿企必须重新评估每一兆瓦电力对应的成本与回报。过去用于挖矿的工业用地、电力基础设施和冷却系统，随着 AI 需求加速增长，可能产生更稳定且更具吸引力的收益。

一旦企业决定重新部署这些资产并进入相关市场，新一轮资本投入便随之开始。销售和使用税、财产税以及用电税等税务因素，可能直接影响转型的投资成本。

本文选取美国部分州，梳理数据中心税收优惠的现状及最新动向，并深入分析得克萨斯州、弗吉尼亚州和佐治亚州等重点辖区的资格门槛与实施路径，为转型中的矿企提供实务参考。

\FloatBarrier
